\subsection*{Correct shapes do not ensure a correct composed answer}
We tested the writer on 571 human answers from DuReader. Of these, 177 answers were contiguous in the supplied question or passage, 94 required rearrangement, and 300 contained a content glyph absent from both inputs. Trained on 2,565 separate questions, a recurrent writer selected a visible span, emitted new ink or stopped. All three decoding conditions produced 25 complete references. Every success was contiguous. Returning the first sentence of the passage without consulting the question gave 66 exact answers. Rotating the questions yielded 24 exact answers with the pixel writer, 13 of them also correct with their original question. The writer missed every exact rearranged and input-absent answer (Fig.~\ref{fig:composed-writer}).

A single blinded reading of twelve questions found one fully supported response among the four input-absent cases for both the pixel writer and copying alone. The pixel writer placed none of 1,525 missing reference cells at the correct output index. With the preceding reference segments supplied and the action forced to write, selection from training shapes recovered 162 of those cells. A learned pixel correction reduced this count to 57. Even with correct preceding segments, the start selector chose an accepted source location on only 399 of 2,634 copy steps and the full span on 235. Both the choice of evidence and the spelling of missing glyphs limit these answers.
